En una zona de l'espai hi ha un camp magnètic uniforme de 0,40~T. En aquesta regió hi ha una espira circular de 200~cm$^2$ d'àrea que gira a 191~rpm (revolucions per minut), tal com indica la figura.

\figura[0.45]{fig1.png}

\begin{apartat}{1}
Si en l'instant inicial el camp magnètic és perpendicular al pla de l'espira, expresseu l'equació del flux magnètic que travessa l'espira en funció del temps.
\end{apartat}

\begin{apartat}{1}
Quina és la força electromotriu (FEM) màxima generada per l'espira?
\end{apartat}
